%%%% EVAL
We report mean compliance, graded binary for four tasks as in the original paper, and strict passes. Our main CoT-Control number is the mean over the seven tasks that are not character-level (all but alternating case, lowercase and uppercase; 350 pairs): across their eight models, \citet{haskins2026obfuscation} find that only these seven predict obfuscation (composite $r=0.898$, against $0.800$ for all ten; $r=0.011$, $-0.197$ and $-0.469$ for the three case tasks). We adopted this composite after the first version of this paper, which used the ten-task mean; every table also reports the ten-task numbers.
%%%% QWEN1
On the seven non-character CoT-Control tasks, mean compliance rises from 13.55 to 19.71, 19.02 and 19.77 (gains of $+6.16$, $[3.73, 8.80]$; $+5.47$, $[2.95, 7.86]$; $+6.22$, $[3.88, 8.55]$), and strict passes from 7 to 63, 63 and 62 of 350; over all ten tasks the mean rises from 15.50 to 24.21, 25.95 and 26.79.
%%%% QWEN2
Within the seven tasks, the seed-42 gain is on third person (84.8\% to 98.2\%), word suppression (4.0\% to 20.0\%), multiple-word suppression (4.0\% to 14.0\%) and repeat sentences (0.0\% to 3.8\%); arrow prefix, end of sentence and meow stay at or near zero as graded. Much of the larger ten-task gain comes from uppercase, a case task (Appendix Table~\ref{tab:standalone-task}).
%%%% PHI
\paragraph{Phi-4-reasoning.} SFT raises ReasonIF IFS from 5.0\% to 12.0\%, almost entirely on the word-limit constraint (without it, $+1.6$, $[-0.4, 3.6]$), and seven-task CoT-Control compliance from 12.40 to 20.66 ($+8.27$, $[5.70, 10.89]$; ten tasks: 9.85 to 16.11), but on the first 100 words no gain remains ($-0.50$, $[-4.56, 4.60]$; Section~\ref{sec:length}), so the gain is explained by shorter outputs.
%%%% OLMO
\paragraph{OLMo-3-7B-Think.} With the Qwen recipe, ReasonIF IFS rises from 8.3\% to 11.7\%, but the seven-task CoT-Control mean does not change ($+0.70$, $[-0.78, 2.28]$; $-1.01$ with degenerate traces scored 0); its ten-task gain ($+2.44$) rested mostly on degenerate passes (14 of its 19 strict passes). With learning rate $2\times10^{-4}$ the style transfers (407 of 500 traces open with a bold heading) and controllability moves in both directions: lowercase rises from 6.2 to 35.5 and uppercase from 4.3 to 16.5, while third person falls from 87.2 to 52.7, plausibly because the training reasoning is written in the first person, so the seven-task mean falls ($-3.01$, $[-5.62, -0.16]$; ten tasks $+1.84$, $[-0.55, 4.32]$; Appendix Tables~\ref{tab:olmo} and~\ref{tab:olmo-task}). On ReasonIF,
%%%% LENGTH
On the seven tasks the Qwen gains are $+4.50$, $+3.94$ and $+2.53$ points for the three seeds, against $+6.16$, $+5.47$ and $+6.22$ over the full traces (ten tasks: $+8.60$, $+9.35$ and $+9.39$ on the first 100 words). This keeps 289 of 350 pairs for seed 42; requiring 200 or 300 words keeps 155 and 73 pairs and gives $+3.81$, $[1.68, 4.34]$ and $+4.18$, $[1.41, 7.49]$.
%%%% LENGTH2
Length therefore explains part of the Qwen gain, but a gain remains for every seed. For Phi the same check removes the gain ($-0.50$ on the seven tasks at 100 words, with 173 of 350 pairs kept).
%%%% DEGEN
 On the seven tasks, scoring degenerate traces 0 gives Qwen gains of $+6.16$, $+5.47$ and $+5.36$ (Table~\ref{tab:models}).
%%%% HEAD1
The CoT-Control grader treats a heading line as a sentence, so a heading can pass or fail a constraint by itself. Grading every trace with a leading heading removed raises the seven-task SFT gains over Base from $+6.16$, $+5.47$ and $+6.22$ to $+10.50$, $+9.21$ and $+9.73$ (ten tasks: $+11.77$, $+13.16$ and $+13.64$), mostly on meow (2.0\% to 22.0\% for seed 42) and end of sentence (0\% to 10.0\%), so the reported gains are conservative; on ReasonIF, removing the heading changes IFS by at most one point. Table~\ref{tab:headings} grades every run this way.

We retrained the reported recipe three more times (one run each). With all 537 headings deleted from the GPT-5.2 data, the seven-task gain disappears ($-1.30$, $[-3.69, 1.18]$); with every heading replaced by one fixed \texttt{**Reasoning**} heading, compliance falls below Base ($-4.60$, $[-7.09, -1.99]$). Both models copy the summaries' first-person opening sentence (``I need to\ldots''), which lowers third person and lowercase below Base. Conversely, placing each prompt's GPT-5.2 opening heading in front of the Qwen3-235B-A22B reasoning (the heading text is copied from the GPT-5.2 example with the same prompt, and the Qwen3-235B-A22B reasoning and answer are left unchanged; nothing is regenerated), which alone gives no gain ($-0.07$, $[-1.44, 1.30]$) and no significant ReasonIF gain, restores part of the effect: $+4.91$, $[2.84, 7.07]$ on the seven tasks against $+10.50$ for the GPT-5.2 summaries, and the full ReasonIF gain (34.7\% IFS, $+20.7$, $[15.0, 26.3]$), with outputs longer than the GPT-5.2 adapter's (mean 3,531 against 2,024 tokens). Over all ten tasks the transplant matches the GPT-5.2 adapter ($+12.57$, $[10.40, 14.85]$), mostly through the case tasks. For Qwen3-235B-A22B, the other-teacher contrast of Table~\ref{tab:variants} is therefore partly one of format rather than teacher. Headings are not the only route, however: Claude Sonnet 4.5 reasoning, which has none, raises seven-task compliance by $+12.13$, $[9.31, 15.01]$ ($+6.70$ with degenerate traces scored 0) while leaving ReasonIF unchanged, and the adapter trained with the loss on the reasoning only writes a heading in every trace but gains little ($+3.03$, $[1.55, 4.55]$; Appendix Table~\ref{tab:seven-all}). On ReasonIF the fixed heading keeps the gain (38.0\%), so there the heading's content matters less than on CoT-Control.

On the case tasks, the headings help because the model fills them in.
%%%% HEAD2
 This account concerns the three case tasks, where the opening's case can be checked directly; they are the tasks that do not predict obfuscation in \citet{haskins2026obfuscation}. On the seven other tasks the heading experiments show that the headings matter, but not how.
%%%% TEACHER
\subsection{Does the Teacher Matter?}
\label{sec:teacher}
The other-teacher runs of Table~\ref{tab:variants} keep the GPT-5.2 prompts but change the reasoning in several ways at once. To test whether a teacher's own controllability carries over, we trained the reported recipe on the same 212 questions with reasoning from two teachers at the extremes of \citet{chen2026cotcontrol}: Claude 3.7 Sonnet (51.1\% stylistic and 49.9\% addition compliance in their Table 2) and DeepSeek R1 (0.2\% and 0.0\%). The reasoning comes from the public s1K releases (\texttt{simplescaling/s1K-claude-3-7-sonnet} and \texttt{simplescaling/s1K-1.1}), which pair the same s1K questions (mostly mathematics and science) with each teacher's raw reasoning, not a summary. We kept the 334 questions whose traces from both teachers fit a 12,288-token training cutoff (the s1K traces are much longer than the summaries) and that share no 8-gram with CoT-Control, and drew 212 at random (seed 42); the cutoff and a batch of $1\times8$ are the only changes to the recipe. ReasonIF is not evaluated, because its questions overlap with the sources of s1K.

The Claude student gains $+7.43$, $[4.89, 10.14]$ on the seven tasks, also on the first 100 words ($+9.07$) and with degenerate traces scored 0 ($+7.72$), whereas the R1 student does not change ($+0.31$, $[-1.05, 1.61]$); paired, the Claude student is $+7.12$, $[4.66, 9.81]$ points above the R1 student (Table~\ref{tab:teacher}). Its gain is on word suppression (4.0\% to 30.0\%) and the arrow prefix (0.0\% to 15.6\%), which the GPT-5.2 summaries do not move, and its traces are longer than Base's (137 of 500 reach the token limit, against 25). Fine-tuning on Qwen3-14B's own reasoning for the GPT-5.2 prompts gives no gain ($+0.98$, $[-0.22, 2.30]$). With the same questions, then, which teacher wrote the reasoning changes the student's controllability, and summaries are not required. For Claude Sonnet 4.5 the comparison is less clear: it is near zero on stylistic and addition constraints in \citet{chen2026cotcontrol} (0.2\% and 0.0\%) but not on suppression (15.6\%, close to Claude 3.7 Sonnet's 15.0\%), and its student's seven-task gain ($+12.13$, $[9.31, 15.01]$) lies mostly on the two suppression tasks and on repeat sentences, many of whose passes are degenerate ($+6.70$ with degenerate traces scored 0); its prompts, length and format also differ from the s1K runs. Both s1K students and the self-distilled model lose most of the base model's lowercase compliance (45.7\% to 16.0\%, 7.8\% and 6.8\%). All of these are single runs on one student model.
%%%% CROSS1
On the seven tasks, with base-derived openings Qwen's continuation compliance is 14.25\% for the base recipient and 13.83\% for the SFT recipient; with SFT-derived openings it is 21.12\% and 18.96\% (ten tasks: 16.11\%, 15.76\%, 29.51\% and 25.26\%; Figure~\ref{fig:haskins-crossed}). Changing the donor from Base to SFT raises the unchanged recipient's compliance by $+6.87$, $[4.37, 9.39]$ points (ten tasks: $+13.40$), while changing the recipient from Base to SFT does not help with the base donor ($-0.41$, $[-1.70, 0.84]$) or the SFT donor ($-2.16$, $[-4.87, 0.70]$; ten tasks: $-0.35$ and $-4.25$). Fifteen SFT$\to$Base and 14 SFT$\to$SFT openings already close the reasoning; excluding them and every unfinished continuation (420 pairs), the ten-task recipient effect is unchanged ($-4.62$, $[-7.43, -1.68]$).
%%%% CROSS2
Prefix-OFF, in which the donor is not given the constraint and its opening names no rule, still reaches 32.0\% IFS ($+18.0$ over Base, $[12.7, 23.3]$; $-1.3$ against SFT, $[-6.3, 3.7]$), and on the seven CoT-Control tasks it keeps most of the donor effect: 19.93\% against 14.25\% after the base donor ($+5.68$, $[3.12, 8.44]$) and 21.44\% for Prefix-ON in the same pipeline ($-1.51$, $[-4.22, 1.31]$). Only over all ten tasks does Prefix-OFF fall well short of Prefix-ON (19.95\% against 29.71\%), through the case tasks. On both benchmarks, then, most of what the SFT opening transfers on the tasks that predict obfuscation does not depend on the donor having seen the constraint.

Phi's opening does not transfer: on the seven tasks the SFT donor changes the base recipient's compliance by $+0.71$, $[-1.15, 2.60]$ points (12.63\% to 13.34\%), whereas the SFT recipient does better than the base recipient with either donor ($+4.46$, $[2.01, 6.90]$; $+8.57$, $[5.85, 11.31]$),
%%%% CTRL1
Over these eight tasks, the transferred ingredient is therefore mainly constraint-specific text at the start of the reasoning, which the SFT donor writes more often.
%%%% CTRL2
These eight tasks include the three case tasks. On the five style tasks that are not character-level (Table~\ref{tab:controls}), the SFT donor raises compliance over the base donor by $+4.42$, $[2.76, 6.21]$ points, but an opening that only names the rule gives nothing ($-1.02$, $[-3.66, 1.57]$ against the base donor), and neither does one that obeys a different rule ($-0.80$, $[-1.80, -0.11]$). The SFT-donor opening written for another question matches the SFT donor ($+0.74$, $[-1.26, 2.79]$), as does the most frequent SFT opening of the task ($-0.80$, $[-2.49, 0.89]$), and only an opening that itself obeys the rule goes further ($+23.41$, $[19.11, 27.84]$ over the SFT donor). On these tasks what the SFT opening transfers is its general style rather than constraint-specific content, consistent with the Prefix-OFF result above.
%%%% DISC
The ablations narrow down why, more clearly on the case tasks than on the seven tasks that predict obfuscation. On the case tasks, compliance is largely decided at the start of the reasoning, in the base model as well as after fine-tuning: training reasoning that begins with a short, free-form span the student can write in the required style (the GPT-5.2 headings, Claude Sonnet 4.5's short opening sentences) raises it, and fixed or ordinary capitalized openings lower it. On the seven other tasks, deleting the headings removes the gain but adding them to another teacher's reasoning restores only part of it, an SFT opening helps whether or not its writer saw the constraint, and the teacher matters: with the same questions, Claude 3.7 Sonnet reasoning raises compliance and DeepSeek R1 reasoning does not. Which property of the training reasoning carries the effect on these tasks remains open.
%%%% LIMIT
 We report the seven-task composite as the main CoT-Control number following \citet{haskins2026obfuscation}, a choice made after seeing the ten-task results; both are in every table. The teacher test uses one pair of teachers on one student, with a longer training cutoff and different questions than the main runs, and the two Claude teachers also differ in prompts, length and format, so it cannot isolate the teacher's controllability.
%%%% AIUSE
 An AI agent also ran the teacher test (training and CoT-Control evaluation of the s1K students), computed the seven-task composite for every run with the repository scripts, and revised the results sections around it.
%%%% INTRO1
on the seven CoT-Control tasks that predict obfuscation in \citet{haskins2026obfuscation}, mean compliance rises from 13.55 to 19.02--19.77 (all ten tasks: 15.50 to 24.21--26.79).
%%%% INTRO2
In Phi-4-reasoning the gain is explained by shorter outputs; in OLMo-3-7B-Think there is no seven-task gain under the same recipe, and at a higher learning rate
%%%% INTRO3
\item \textbf{The headings matter} (Section~\ref{sec:headings}). Deleting the summaries' bold headings from the training data removes the CoT-Control gain and a fixed heading lowers compliance below Base; adding the GPT-5.2 headings to Qwen3-235B-A22B reasoning reproduces the ReasonIF gain and part of the seven-task gain. On the case tasks, compliance across twelve Qwen3-14B runs tracks how often the model's first heading or sentence is itself written in the required case.
\item \textbf{The teacher matters} (Section~\ref{sec:teacher}). On the same 212 questions, raw reasoning from Claude 3.7 Sonnet, among the most controllable models in \citet{chen2026cotcontrol}, raises seven-task compliance by $+7.43$ points, and reasoning from DeepSeek R1, among the least, does not; summaries are not required.
%%%% INTRO4
An opening written by the SFT model raises the unchanged base model's seven-task continuation compliance from 14.25\% to 21.12\%, while swapping in the fine-tuned weights does not help. On these tasks the opening helps whether or not its writer saw the constraint, and an SFT opening written for another question works as well; constraint-specific opening text matters mainly on the case tasks.
%%%% SELFHEAD
 Placing the same headings in front of Qwen3-14B's own reasoning for the GPT-5.2 prompts, which alone gives no gain ($+0.98$, $[-0.22, 2.30]$), gives a smaller one ($+3.92$, $[1.75, 6.30]$ on the seven tasks; ReasonIF IFS 28.7\%, $+14.7$, $[9.3, 20.3]$) and restores most of the base model's lowercase compliance (34.8\% against 6.8\% without the headings and 45.7\% for Base).
%%%% PHITEACH
On Phi-4-reasoning the Claude 3.7 student also gains more than the R1 student ($+8.56$ against $+3.41$; paired, $+5.15$, $[2.76, 7.74]$), but, as for Phi's GPT-5.2 gain (Section~\ref{sec:length}), the difference disappears on the first 100 words ($+1.23$, $[-2.63, 3.40]$). All of these are single runs.
%%%% RIFOV
ReasonIF is not evaluated, because the two question sets overlap: 4 of the 300 ReasonIF questions appear verbatim among the s1K questions (1 among the 212 used) and 37 share an 8-gram with one (24 with the 212), mostly AIME and AMC problems.
%%%% OPEN
Across all 16 Qwen3-14B runs (Base, three seeds, the four heading runs, the five Table~\ref{tab:variants} adapters, self-distillation and the two s1K students of Section~\ref{sec:teacher}), the number of lowercase and uppercase items whose opening, the first heading or else the first sentence, is written in the required case predicts the run's compliance on that task (Spearman $\rho = 0.91$ and $0.82$), and grading only the reasoning after the opening, compliance is 82.5\% after an obeying opening against 22.0\% after another one. The exception is self-distillation with the GPT-5.2 headings, which writes 46 of its 50 lowercase openings in lowercase but complies with only 37.7\% after them.
%%%% APPRUNS
 For the teacher test (Section~\ref{sec:teacher}), Qwen3-14B and Phi-4-reasoning were trained with the same recipe on the two s1K sets, with a cutoff of 12,288 tokens and a per-device batch of 1 with 8 accumulation steps (the same effective batch of 8). No example exceeds the cutoff under the Qwen3-14B tokenizer, by construction of the selection; under Phi-4-reasoning's tokenizer 1 DeepSeek R1 example exceeds it by 157 tokens, which truncates the end of its answer, and 0 Claude 3.7 Sonnet examples do. The self-distillation adapters use Qwen3-14B's own reasoning and answers for the 212 GPT-5.2 prompts.
%%%% INTROOPEN
On the case tasks, compliance across 16 Qwen3-14B runs tracks
%%%% ABSTRACT
Chain-of-thought (CoT) controllability, a model's ability to follow instructions about its own reasoning, matters for CoT monitoring because it correlates with the ability to obfuscate reasoning. We ask whether ordinary supervised fine-tuning (SFT) on benign reasoning data, with no controllability objective, can raise it. We fine-tune Qwen3-14B, Phi-4-reasoning, and OLMo-3-7B-Think on 212 GPT-5.2 reasoning summaries of software-engineering requests and evaluate them on ReasonIF and on the seven CoT-Control tasks that predict obfuscation. In Qwen3-14B, ReasonIF instruction-following success rises from 14.0\% to 33.3\% (33.7\% and 33.0\% for two further training seeds), and CoT-Control compliance rises from 13.55 to 19.02--19.77; the CoT-Control gain survives scoring only the first 100 words and scoring degenerate outputs as failures, and ReasonIF accuracy falls from 80.3\% to 59.3--65.3\%. Ten tokens written by the fine-tuned model raise the unchanged base model's CoT-Control compliance by 6.87 points. The teacher matters: on another 212 questions, raw reasoning from Claude 3.7 Sonnet raises CoT-Control compliance by 7.43 points, while reasoning from DeepSeek R1 does not. In Phi-4-reasoning the gain is explained by shorter outputs, and OLMo-3-7B-Think shows none. Ordinary distillation can therefore raise controllability, which should be re-measured after post-training.
%%%% CONCL
Supervised fine-tuning on 212 benign GPT-5.2 reasoning summaries raised measured CoT controllability in Qwen3-14B, on ReasonIF and on the seven CoT-Control tasks that predict obfuscation, consistently across three training seeds and beyond what shorter outputs or degenerate text explain, and ten tokens written by the fine-tuned model transfer much of the CoT-Control gain to the unchanged base model. Which teacher wrote the training reasoning matters: with the same questions, raw reasoning from Claude 3.7 Sonnet raises controllability and reasoning from DeepSeek R1 does not. The summaries' opening headings matter as well, but they account for the case tasks better than for the tasks that predict obfuscation, where what carries the effect remains open. The effect is model-dependent: it is explained by shorter outputs in Phi-4-reasoning and absent in OLMo-3-7B-Think. Controllability should be re-measured after ordinary distillation, not assumed to fall.
